\documentclass[aspectratio=169,10pt]{beamer}
\usetheme{Madrid}
\usecolortheme{default}

\usepackage{hyperref}
\usepackage{tikz}
\usetikzlibrary{arrows.meta,decorations.pathmorphing,backgrounds,positioning,fit,petri}
\usepackage{amsmath,amssymb}

% --- Custom Colors ---
\definecolor{unibblue}{RGB}{0,51,102}
\definecolor{unibgold}{RGB}{212,175,55}
\definecolor{unibgreen}{RGB}{0,128,0}

\setbeamercolor{palette primary}{bg=unibblue,fg=white}
\setbeamercolor{palette secondary}{bg=unibblue!85,fg=white}
\setbeamercolor{palette tertiary}{bg=unibblue!70,fg=white}
\setbeamercolor{structure}{fg=unibblue}
\setbeamercolor{block title}{bg=unibblue!85,fg=white}
\setbeamercolor{block body}{bg=unibblue!10,fg=black}
\setbeamercolor{alertblock title}{bg=unibgold!90,fg=black}
\setbeamercolor{alertblock body}{bg=unibgold!15,fg=black}

\title{Persamaan Diferensial}
\subtitle{Minggu 3: Aplikasi PDB Orde 1 dalam Rekayasa Elektro (Transien RC, RL, \& Termal)}
\author{Novalio Daratha \& Muhammad Arfan}
\institute{Program Studi Teknik Elektro\\Universitas Bengkulu}
\date[\today]{2026 \\ \vspace{2mm}\small\textit{Tanggal Kompilasi: \today}}

\begin{document}

\begin{frame}
  \titlepage
\end{frame}

\begin{frame}{Capaian Pembelajaran Perkuliahan (CPMK 2)}
  \begin{itemize}
    \item \textbf{CPMK 2.1:} Mahasiswa mampu menyusun pemodelan fisis rangkaian elektrik orde 1 menggunakan Hukum Kirchhoff (KVL/KCL).
    \item \textbf{CPMK 2.2:} Mahasiswa mampu menyelesaikan PDB orde 1 untuk respons transien dan kondisi mantap (\textit{steady-state}).
    \item \textbf{CPMK 2.3:} Mahasiswa memahami makna fisis dari konstanta waktu (\textit{time constant} $\tau$) pada pengisian/pengosongan energi.
    \item \textbf{CPMK 2.4:} Mahasiswa mampu memverifikasi hasil analitik dengan simulasi numerik berbasis bahasa Julia.
  \end{itemize}
\end{frame}

\begin{frame}{Daftar Isi Materi Minggu 3}
  \tableofcontents
\end{frame}

\section{Pemodelan Fisis Rangkaian RL Seri}
\begin{frame}{Rangkaian RL Seri: Pengisian Arus Induktor}
  \begin{columns}
    \begin{column}{0.5\textwidth}
      \begin{block}{Persamaan KVL}
        Berdasarkan Hukum Tegangan Kirchhoff:
        $$ V_R(t) + V_L(t) = V_0 $$
        $$ R \cdot i(t) + L \frac{di(t)}{dt} = V_0 $$
      \end{block}
      \pause
      Bentuk standar PDB linear:
      $$ \frac{di}{dt} + \frac{R}{L}i = \frac{V_0}{L} $$
    \end{column}
    \begin{column}{0.5\textwidth}
      \begin{alertblock}{Karakteristik Fisik}
        \begin{itemize}
          \item Induktor menentang perubahan arus tiba-tiba: $i(0^-) = i(0^+) = 0$.
          \item Konstanta waktu: $\tau = \frac{L}{R}$ (detik).
          \item Arus akhir saat mantap ($t \to \infty$): $I_{\text{maks}} = \frac{V_0}{R}$.
        \end{itemize}
      \end{alertblock}
    \end{column}
  \end{columns}
\end{frame}

\begin{frame}{Penurunan Solusi Analitik Arus Induktor}
  Menggunakan faktor integrasi $\mu(t) = e^{\int \frac{R}{L} dt} = e^{\frac{R}{L}t}$:
  \pause
  \begin{equation*}
    e^{\frac{R}{L}t}\frac{di}{dt} + \frac{R}{L}e^{\frac{R}{L}t}i = \frac{V_0}{L}e^{\frac{R}{L}t} \implies \frac{d}{dt}\left(i \cdot e^{\frac{R}{L}t}\right) = \frac{V_0}{L}e^{\frac{R}{L}t}
  \end{equation*}
  \pause
  Integralkan kedua ruas terhadap $t$:
  \begin{equation*}
    i(t) \cdot e^{\frac{R}{L}t} = \int \frac{V_0}{L}e^{\frac{R}{L}t} dt = \frac{V_0}{L}\left(\frac{L}{R}\right)e^{\frac{R}{L}t} + C = \frac{V_0}{R}e^{\frac{R}{L}t} + C
  \end{equation*}
  \pause
  \begin{block}{Solusi Umum \& Solusi Khusus IVP ($i(0)=0$)}
    $$ i(t) = \frac{V_0}{R} + C e^{-\frac{R}{L}t} \implies 0 = \frac{V_0}{R} + C \implies C = -\frac{V_0}{R} $$
    $$ {\color{unibblue} i(t) = \frac{V_0}{R}\left(1 - e^{-t/\tau}\right)}, \quad \tau = \frac{L}{R} $$
  \end{block}
\end{frame}

\section{Pemodelan Fisis Rangkaian RC Seri}
\begin{frame}{Rangkaian RC Seri: Pengisian Muatan \& Tegangan}
  \begin{columns}
    \begin{column}{0.5\textwidth}
      \begin{block}{Persamaan Diferensial}
        $$ V_R(t) + V_C(t) = V_0 $$
        $$ R \frac{dq}{dt} + \frac{q}{C} = V_0 \implies R C \frac{dV_C}{dt} + V_C = V_0 $$
      \end{block}
      \pause
      Konstanta waktu kapasitor:
      $$ \tau = R \cdot C \quad (\text{detik}) $$
    \end{column}
    \begin{column}{0.5\textwidth}
      \pause
      \begin{alertblock}{Solusi Pengisian ($V_C(0) = 0$)}
        $$ V_C(t) = V_0 \left(1 - e^{-t/RC}\right) $$
        $$ i(t) = C \frac{dV_C}{dt} = \frac{V_0}{R} e^{-t/RC} $$
      \end{alertblock}
      \textit{Tegangan naik mulus, arus melonjak maksimum lalu meluruh.}
    \end{column}
  \end{columns}
\end{frame}

\begin{frame}{Pengosongan Kapasitor (Discharging Transient)}
  Bila kapasitor yang telah terisi penuh $V_0$ dihubungkan singkat melalui resistor $R$ ($V_{\text{sumber}} = 0$):
  $$ R C \frac{dV_C}{dt} + V_C = 0 \implies \frac{dV_C}{V_C} = -\frac{1}{RC} dt $$
  \pause
  Integrasikan dengan kondisi awal $V_C(0) = V_0$:
  \begin{block}{Respons Pelepasan Energi}
    $$ \ln\left(\frac{V_C(t)}{V_0}\right) = -\frac{t}{RC} \implies {\color{unibblue} V_C(t) = V_0 e^{-t/\tau}} $$
  \end{block}
  \pause
  \begin{itemize}
    \item Saat $t = 1\tau$: $V_C(t) = V_0 e^{-1} \approx 0.368 V_0$ (Tersisa $36.8\%$).
    \item Saat $t = 5\tau$: $V_C(t) = V_0 e^{-5} \approx 0.0067 V_0$ (Sistem dianggap mencapai \textit{steady-state} $<1\%$).
  \end{itemize}
\end{frame}

\section{Aplikasi Termal: Hukum Pendinginan Newton}
\begin{frame}{Hukum Pendinginan Newton pada Trafo / Motor Listrik}
  \begin{block}{Persamaan Laju Pelepasan Kalor}
    Laju perubahan suhu $T(t)$ konduktor terhadap suhu sekitar $T_a$:
    $$ \frac{dT}{dt} = -k(T - T_a) $$
  \end{block}
  \pause
  Penyelesaian menggunakan separasi variabel:
  $$ \int \frac{dT}{T - T_a} = -k \int dt \implies \ln|T - T_a| = -kt + C_1 $$
  \pause
  \begin{alertblock}{Solusi Suhu Motor/Kabel Terhadap Waktu}
    Jika suhu awal konduktor adalah $T(0) = T_0$:
    $$ T(t) = T_a + (T_0 - T_a)e^{-kt} $$
  \end{alertblock}
\end{frame}

\section{Uji Konsep \& Kuis Interaktif}
\begin{frame}{Kuis Interaktif 1: Konstanta Waktu Rangkaian}
  \textbf{Pertanyaan:} Sebuah rangkaian pengapian mobil menggunakan induktor $L = 500\text{ mH}$ dan resistor $R = 25\ \Omega$. Berapakah konstanta waktu $\tau$ dan berapa waktu yang diperlukan arus untuk mencapai $99\%$ dari nilai maksimumnya?
  \pause
  \begin{block}{Pembahasan:}
    \begin{enumerate}
      \item $\tau = \frac{L}{R} = \frac{0.5\text{ H}}{25\ \Omega} = 0.02\text{ detik} = 20\text{ ms}$. \pause
      \item Arus mencapai $99\%$ saat: $1 - e^{-t/\tau} = 0.99 \implies e^{-t/\tau} = 0.01$. \pause
      \item $t = -\tau \ln(0.01) \approx 4.605 \times 20\text{ ms} \approx 92.1\text{ ms} \approx 5\tau$.
    \end{enumerate}
  \end{block}
\end{frame}

\begin{frame}{Integrasi Komputasi: Praktikum Julia}
  \begin{block}{Jupyter Notebook Minggu 3: Simulasi Transien RC \& RL}
    \begin{itemize}
      \item Mempelajari algoritma \textbf{Metode Euler} secara langsung:
      $$ i_{k+1} = i_k + \Delta t \cdot \left(\frac{V_0 - R i_k}{L}\right) $$
      \item Membandingkan galat (*truncation error*) antara solusi numerik dan analitik eksak.
      \item Menjalankan pustaka industri \textbf{DifferentialEquations.jl} (SciML).
    \end{itemize}
  \end{block}
\end{frame}

\begin{frame}{Ringkasan \& Kesimpulan}
  \begin{enumerate}
    \item Rangkaian RL dan RC dimodelkan oleh PDB linear orde 1 dengan konstanta waktu $\tau = L/R$ atau $\tau = RC$.
    \item Respons total selalu terdiri atas **respons alami (transien)** yang meluruh dan **respons paksa (keadaan mantap)**.
    \item Mahasiswa wajib memverifikasi perhitungan matematis menggunakan simulasi Julia pada Lembar Kerja Mandiri.
  \end{enumerate}
\end{frame}

\end{document}
